\section{Introduction}
\label{2406:sec:main:introduction}

Descent-based algorithms, such as Stochastic Gradient Descent (SGD) and its variations, are the backbone of contemporary machine learning. Their simplicity in implementation, efficiency in operation, and notably effective performance in practice highlight their importance. A mathematical understanding of SGD's effectiveness remains a key focus in the field. Recent progress has been particularly noteworthy in the realm of shallow neural networks. A sequence of works demonstrated that optimizing large width two-layer neural networks can be mapped into a convex optimization problem over the space of probability measures of weights, the so-called mean-field analysis \citep{mei2018mean,chizat2018global,rotskoff2022trainability,sirignano2020mean}. 
Following this breakthrough, a large part of the theoretical effort has shifted to describing what class of functions can be efficiently learned by SGD, i.e. time and computational complexities required to learn a given class of functions. This has been, in particular, 
thoroughly analyzed in a series of recent works focusing on isotropic distributions (e.g. Gaussian, spherical or in the hypercube) and targets depending only on a few relevant directions (a.k.a. \emph{multi-index models}). A key result from this literature is that the time complexity of SGD scales with the covariates dimension according to the so-called \emph{information exponent} \citep{arous2021online} for single-index and \emph{leap complexity} \citep{abbe2021staircase,abbe2023sgd} for multi-index targets, sparking increasing interest from the theoretical machine learning community over the last few months \citep{damian2022neural,damian2023smoothing,dandi2023how,bietti2023learning,ba2023learning,moniri2023theory,mousavihosseini2023gradient,zweig2023symmetric}. \\
Our work follows this thread, focusing instead on the effect of batch size $n_b$, parallelization, and sample-splitting into the overall complexity required to learn a multi-index target. Instead of looking at data one-by-one, as is common in theoretical studies, we investigate the finite $n_b$ problem, and characterize the time/complexity tradeoff when learning with one-pass SGD. 
% with large $n_b = O(d)$ minimizes the training time without changing the total sample complexity.  
Our central goal is to paint a complete picture 
of how fast generalized linear models and two-layer neural networks adapt to the features of training data as a function of $n_b$, and the structure of the target function. 

Our analysis sheds light on a fundamental limitation of one-pass (or online) SGD, namely that for batch sizes larger than the input dimension, the dynamics of the training algorithm is dominated by negative feedback terms that do not permit to reduce the time iterations needed to learn the target. Therefore, we provide a rigorous solution to this fundamental limitation of SGD by considering gradient updates on the correlation loss. 
% This assumption has been used in similar context \cite{damian_2023_smoothing}, and we refer to this modified training protocol as \emph{Correlation Loss SGD}. 
Our approach, drawing inspiration from the \textit{summary statistics} method employed by \cite{saad1995line,arous2021online,arous2022high}, concentrates on the overlaps of neurons with the target subspace and their norms. This differs from recent studies, such as those by \cite{abbe2022merged} and \cite{damian2022neural}, which focus on the full gradient vector.

% Removing the self-correlation term from the loss is often presented as a necessary mathematical assumptions for the analysis.

% We provide precise characterizing of the benefits of including it or not for different scaling of the batch size \(n_b\) and the learning rate \(\gamma\).
